\documentclass{article}
\usepackage{graphicx} % Required for inserting images
\usepackage{amssymb}

\title{Fundamentals Chapter 1}
\author{Siddharth K}
\date{September 2024}
\newcommand{\FF}{\mathbb{F}}
\newcommand{\ZZ}{\mathbb{Z}}
\newcommand{\QQ}{\mathbb{Q}}
\newcommand{\RR}{\mathbb{R}}
\newcommand{\CC}{\mathbb{C}}

\begin{document}

\maketitle

\textbf{1a.} Since $1 + (-1) = 0$, $0a = (1 + (-1))a = 1\cdot a + (-1)\cdot a = a + (-a) = 0$. \\

\textbf{1b.} Assume for contradiction that $0 = 1$. For all $a \in \FF$, $a\cdot 0 = a\cdot 1$. We know that $a\cdot 0 = 0$ and $a\cdot 1 = a$ by the definition of $1$, so $a = 0$. But this implies that the only element in $\FF$ is $0$, which contradicts the definition of a field. Therefore, $0\neq 1$. \\

\textbf{1c.} Assume $ab = 0$ and $a, b \neq 0$ so that $a^{-1}$ and $b^{-1}$ exist. Then $aba^{-1}b^{-1} = 0a^{-1}b^{-1}$, so $1 = 0$, which is a contradiction. Therefore at least one of $a, b$ must be $0$. \\

\textbf{1d.} We know that $0a = 0b$. Replacing $0$ with $-a + a$ and $-b + b$ gives us $$a(-b + b) = b(-a + a),$$ $$a(-b) + ab = (-a)b + ab,$$ $$a(-b) + ab + (-ab) = (-a)b + ab + (-ab),$$ $$a(-b) = (-a)b.$$
Also, from $0a = 0$, we can say that $$0a = 0,$$ $$a(b + (-b)) = 0,$$ $$ab + a(-b),$$ $$ab + (-ab) + a(-b) = -ab,$$ $$a(-b) = -ab.$$ Therefore, $a(-b) = (-a)b = ab$. \\

\textbf{2.} As shown earlier, $0x = 0$ for all $x \in \FF$, and $0 \neq 1$. Therefore, $x \cdot 0 = 0 \neq 1$. \\

\textbf{3.} This is equivalent to the statement that $a\cdot b = 0$ implies $a$ or $b = 0$ iff for all $a \in S$ there exists $a^{-1}$ such that $a\cdot a^{-1} = 1$. \\

First, assume that $a^{-1}$ exists for all $a \in S$. Now, assume $ab = 0$, and assume $a, b \neq 0$. Since they are both nonzero, $a^{-1}$ and $b^{-1}$ exist. This gives us $aba^{-1}b^{-1} = 0a^{-1}b^{-1}$, so $1 = 0$, which is a contradiction. Therefore at least one of $a, b$ must be $0$. \\

Now, assume that if $ab = 0$ then $a$ or $b = 0$. Then assume $a_1 b = a_2 b = c$. Then $$a_1 b + (-a_2b) = 0,$$ $$(a_1 + (-a_2))b = 0.$$ Because of our initial assumption, either $a_1 = a_2$ or $b = 0$. This implies that either $b = 0$, or if $ab = c$ for some fixed $c$ then $a$ is unique. \\

Now, say $S$ has $n$ elements and call them $a_1, a_2, \ldots, a_n$. Multiplying them by $b$ gives us $a_1b, a_2b, \ldots, a_nb$. Since each $a_i$ is unique, each $a_ib$ is unique as well. Since there are $n$ values of $a_ib$ and there are $n$ elements in $S$, exactly one of them is $1$. The $a_i$ such that $a_ib = 1$ is the inverse of $b$. \\

\textbf{6.} Assume $\FF$ is a subfield of $\QQ$. Then it must satisfy all field axioms with the same operations of $\QQ$. \\

Since $\FF$ is a field, it must contain the additive identity, $0$, and the multiplicative identity, $1$ in $\QQ$. By repeatedly adding $1$, we obtain all positive integers. Also, $1$ has an additive inverse, $-1$. By repeatedly adding $-1$, we obtain all negative integers. Therefore, $\FF$ contains $\ZZ$. \\

Also, $\FF$ must contain the multiplicative inverses of all non-zero elements it contains, so if $n \in \FF$ then $\frac{1}{n} \in \FF$. So $\FF$ contains all fractions $\frac{a}{b}$ where $a, b \in \ZZ$ and $b \neq 0$, which are all the elements of $\QQ$. Since $\FF$ contains all elements of $\QQ$, it follows that $\FF = \QQ$, and there is no proper nontrivial subfield of $\QQ$. \\

\textbf{7.} Assume, for contradiction, that there is a field $\FF$ satisfying $\RR \subsetneq \FF \subsetneq \CC$. So there is an element $z \in \FF$ so that $z \not \in \RR$ but $z \in \CC$. This means $z$ can be written as $z = a + bi$ where $a, b \in \RR$ and $b$ is nonzero. \\

Isolating $i$ gives us $i = \frac{z - a}{b}$. Since $a, b \in \RR$, $a$ has an additive inverse and $b$ has a multiplicative inverse. So $\frac{z - a}{b}$ exists in $\FF$, so $i \in \FF$. \\

Since $\FF$ contains all elements of $\RR$ and it contains $i$, it contains all elements that can be obtained by adding $ai$ or $-ai$ to any real number where $a \in \RR$. This gives us all numbers of the form $c_1 + c_2i$, which is all of $\CC$. Therefore, $\FF = \CC$, which is a contradiction. So there is no field $\FF$ satisfying $\RR \subsetneq \FF \subsetneq \CC$. \\

\textbf{8.} We need to prove that the axioms of a field apply to $\ZZ/p\ZZ$. All the axioms except (M2) hold because they are properties of integers. \\

From problem 1.3, if a finite set along with two operations satisfy everything except (M2) and it is true that $ab = 0$ implies that $a$ or $b = 0$. So it suffices to show that if $ab \equiv 0 \pmod{p}$ then either $a$ or $b \equiv 0 \pmod{p}$. Assume that $a, b \not \equiv 0 \pmod{p}$. Then $ab \not \equiv 0 \pmod{p}$, which is a contradiction. So $\ZZ/p\ZZ$ is a field. \\

\textbf{11.} Since $G$ is a group, it satisfies $(gh)k = g(hk)$. So if $g(hk) = (gh)(gk)$ then $(gh)k = (gh)(gk)$. Multiplying by the inverse of $gh$ gives us $k = gk$ for all $g, k \in G$. This means $g$ is the identity element. Since this is true for all $g \in G$, $G = \{e\}$. \\

\textbf{12.} Assume axioms (A1R) and (A1L) to be true. Then $g\cdot e_R = g$. If we let $g = e_L$, we get $e_L\cdot e_R = e_R$. Similarly, letting $g = e_R$ in $e_L\cdot g = g$ gives us $e_L\cdot e_R = e_L$. Therefore, $e_R = e_L\cdot e_R = e_L$, so $e_R = e_L$. This means (A1R) and (A1L) imply (A1). \\

To go the other direction, assume there exists $e \in G$ so that for all $g \in G$, $g\cdot e = e\cdot g = g$. Then letting $e_R = e_L = e$ gives us $g\cdot e_L = e_R\cdot g = g$ which implies group axioms (A1R) and (A1L).

\end{document}